% !TeX root = ../papers/nc-p3.tex
\section{滑らかな二次曲面商を持つ全窓と二次元射影像の次数}
\label{sec:smooth-quadric-quotients}
滑らかな二次曲面のグラフが点対応の一成分でよく、余分な点、直線、非被約構造を許す。
不変な仮定は、全整数位置に、滑らかな二次曲面の斉次座標環 $D$ の反復への
一次で同型な全射が存在することである。グラフの移動は超平面束を保つので
互換なグラフを含む場合は位置別基底を運べばこの形となる。
\begin{theorem}\label{thm:smooth-quadric-quotients}
この商を持つ標準14641型の全添字保持同型類は以下で尽き、全て最小周期1である。
\begin{enumerate}
\item 二つの可逆二次行列 $P,Q$ が同じ特性多項式を持つ場合。各々の共役、
両者の交換と共通スカラーだけを除く。相異なる二根なら根の比
$u\in k^\times\setminus\{1\}$ を $u\sim u^{-1}$ で割ったもの。
重根なら両方スカラー、片方だけJordan型、両方Jordan型の三点である。
\item $P=\operatorname{diag}(1,-1)$, $Q=\operatorname{diag}(r,-r)$,
$r\in k^\times$ とし、次位置で $r$ を $r^{-1}$ に替えて交互に繰り返す場合。
$r\sim-r\sim r^{-1}$ だけを除く。この枝の $r^2=1$ は第一の $u=-1$ と同じ一点であり、
それ以外では前節の二次曲面だけのグラフとなる。
\end{enumerate}
表示からの復元では、下記の七次元空間中の各 $(P_i,Q_i)$ を二次双対の追加関係に取る。
両方スカラーの場合は可換多項式窓であり、それ以外では二次曲面は点対応の唯一の二次元既約成分として復元できる。
この商を持つ有限Frobenius候補には非Koszul例はない。
\end{theorem}
\begin{proof}[四次対合による可逆性の強制]
$D^!$ の通常のClifford表示 $C$ と中心 $h$ は前節のものを使う。
$h=1$ での評価は各斉次片で単射である。実際、$k[h]$ 上の十六基底では
各語の係数は同じ次数の $h$ の冪の定数倍しか取り得ない。
$V^*$ の二次形式を $2\times2$ 行列の行列式とする。一次片のブロック表示は
\[
 v(U)=\begin{pmatrix}0&U\\ U^\#&0\end{pmatrix},\qquad
 U^\#=\operatorname{adj}(U),\qquad U\in M_2(k).
\]
この行列の平方は $\det(U)$ 倍の単位なのでClifford関係を満たす。
二つの行列単位の積で偶ブロックを取り、さらに一次を掛ければ全行列環が生成される。
具体的に $v(U)v(V)$ の対角ブロックは $(UV^\#,U^\#V)$ で跡が等しく、その線型包は
\[
 H=\{(P,Q)\in M_2(k)\oplus M_2(k):\operatorname{tr}P=\operatorname{tr}Q\}.
\]
行列単位で確認するか、$C_2,C_3,C_4$ の次元 $7,8,8$ と四次での体積を使ってもよい。
$C_3$ は非対角の全二ブロック、$C_4$ は対角の全二ブロックとなる。
$K$ を曲面のグラフを消す七つの関係の空間とすれば、窓の二次双対は
$C$ の反復で各位置の一つの非零 $a_i=(P_i,Q_i)\in H$ を消したものである。

四次の非零traceを $C_4$ に引き戻し、
$\operatorname{str}(b_i\,\cdot)$ と表す。ここで $b_i$ は二ブロック行列、
$\operatorname{str}(M,N)=\operatorname{tr}M-\operatorname{tr}N$ である。
$H^\perp=kI$ なので、$H\times H$ 上の四次積対合の左核は
$\{a\in H:b_i a\in kI\}$ である。標準Frobenius性からこの核は正確に $ka_i$、
右核は $ka_{i+2}$ である。$b_i$ が可逆なら $a_i$ も可逆である。
$b_i$ が特異なら $b_i a\in kI$ は $b_i a=0$ に限る。二ブロックの階数の和が
2以下なら全偶空間内の核は4次元以上で、超平面 $H$ 内でも3次元以上となり得ない。
従って階数は1と2である。その核内の $a_i$ は一ブロックだけに台を持ち、
その像は直線、跡は零の非零行列である。特に、どの $b_i$ も両ブロックが非零である。

しかしそのような一ブロック $a_i$ は前の四次traceと矛盾する。
$C_1a_i C_1$ は、$U,V$ が任意の二行二列行列であることから、
反対側の一ブロック全体を張る。任意の非零行列 $P$ で $M_2 P M_2=M_2$ だからである。
前のtraceはこの全空間を消すので $b_{i-1}$ の一ブロックが零になってしまう。
従って全 $b_i$ と全 $a_i$ は可逆である。この議論では各点への評価階数を仮定していない。
\end{proof}

\begin{proof}[二ブロックの移送と全特殊値]
三次関係イデアルの次元は $8-4=4$ であり、可逆性により $a_i C_1$ と $C_1 a_{i+1}$ は
ともに四次元で等しい。従って任意の $U$ で $a_i v(U)a_{i+1}^{-1}$ の
下ブロックが上ブロックの余因子となる。式は
\[
 Q_i U^\# P_{i+1}^{-1}
 =(Q_{i+1}^{-1})^\# U^\# P_i^\# .
\]
全 $U^\#$ に対する左右乗法の一致から $[P_{i+1}]=[P_i]$,
$[Q_{i+1}]=[Q_i]$ となる。$a_i$ のスカラーを除いて
$a_i=(P,\lambda_i Q)$ と書くと
\[
 \lambda_i\lambda_{i+1}=\det P/\det Q,\qquad
 \operatorname{tr}P=\lambda_i\operatorname{tr}Q.
\]
跡が非零なら倍率は一定で、下ブロックを取り直せば二行列の跡も行列式も等しい。
ともに跡零なら可逆だから異なる二根の対角形に取れ、表示の $r,r^{-1}$ が得られる。
初めから一定で跡零の場合はちょうど $r^2=1$ である。

逆にこのリストの各解では、移送式が成立し $a_i$ は各斉次片で正則な正規関係となる。
可逆性は $h=1$ の行列と各斉次片の単射性から読める。さらに両ブロックの逆行列は跡が等しい。
跡零の場合は両方零、それ以外の場合は行列式も等しいからである。
逆対を $C_2$ の元 $b$ へ持ち上げると $a_i b=b a_i=h^2$ である。
従って前節で用いた $C/(h^2)$ のFrobenius対合への埋込みが同様に働き、
商双対は級数 $(1+z)^4$ のKoszulかつFrobeniusである。
同じ次数の正規二次関係によるTor二段の議論はそのまま使える。これで全次数の正則性を得る。

任意の非スカラー表示では $\Omega_i\ne0$ である。階数4なら辺はちょうど曲面のグラフ。
階数2なら $p\wedge y=t\Omega_i$ の $t\ne0$ の解は、$\Omega_i$ の台である同じ直線に
両点が乗る場合に限る。そこでは $q(p,y)=t c_i$ の双一次式が残る。
この式が直線の積全体で消えるには $c_i=0$ かつ台が全等方でなければならない。
この場合 $a_i$ は直交する二つの等方文字のClifford積の非零倍となり平方零なので、可逆性と矛盾する。
従って余分な部分の次元は高々1で、最大の二次元成分から $Q$ とそのグラフが一意に戻る。
第二列の $r^2\ne1$ が曲面だけであることは前節の $U+mT$ での両ブロックの固有値が
それぞれ $\pm(1+m),\pm(1-m)$ であることによる。
同型は各位置で同じ曲面の射影変換である。
連結部分は $U\mapsto G U H^{-1}$（$G,H\in\mathrm{SL}_2$）で両ブロックに独立な共役を誘導し、
$U\mapsto U^\#$ がブロックを交換する。従ってリストにある同一視でちょうど尽きる。
共通スカラーだけが正規化されていないことに注意する。両方スカラーのときは追加関係が $h$ で
元の窓は可換多項式表示である。一移動は一定列なら恒等、跡零交互列ならブロック交換まで恒等である。
\end{proof}

\begin{proposition}\label{prop:surface-degree-bound}
四辺の各射影のファイバーが空または一点なら、前節と同様に全辺の共有射影像は
スキームとして一致し、点対応は相互の同型のグラフとなる。
一辺の第一射影像が平面三空間全体でなければ、その二次元部分の超曲面因子の総次数は
重複込みで高々2である。特に滑らかな二次曲面が一つの既約成分として現れる標準窓は、
余分な埋込み成分を許しても上のリストで尽くされる。
\end{proposition}
\begin{proof}
共有像の一致には双対の標準縮約を用いる。四次元文字に作用する六関係の評価行列を
$M(p)$ ($6\times4$) とし、一般階数4とする。その四次小行列は四次式である。
共通四次因子があれば四平面 $\operatorname{im}M(p)$ のPlücker点が一定となる。
よって六関係に定数の線型依存があり矛盾する。
共通三次因子があれば、それを除いたPlücker座標は一次式で、像の閉包は
$\operatorname{Gr}(4,6)$ 内の射影線型空間となる。その四平面全てが固定五平面内か、
全てが固定三平面を含む。実際二点を結ぶ直線がGrassmann多様体に入るには平面対が
三次元で交わる必要があり、二平面の共通三平面を含まない第三平面はその和五平面に入る。
両型を同時に外へ広げると三次元交叉条件に反する。固定五平面の場合は再び六関係が従属である。

固定三平面で割ると $3\times4$ の一次式行列の全評価が階数1以下である。
階数一行列の線型空間では像直線が共通か行直線が共通である。前の場合は既に除外した。
後の場合、三つの独立な関係が $\ell_j(p)v(y)$ と同じ右因子を持つ。
平面 $v(y)=0$ 上、残る三関係が左文字に作用する $3\times4$ 行列は
各点で階数3を持たなければならない。そうであれば
$\mathcal O^4\twoheadrightarrow\mathcal O(1)^3$ の核は $\mathcal O(-3)$ である。
平面上で全Chern類を掛けると $(1-3H)(1+H)^3=1-6H^2\ne1$ で矛盾する。
従って三次因子もない。
滑らかな二次曲面の成分を次の位置に送った像は、同型な二次元の非平面的整曲面で、
上の次数上界から再び滑らかな二次曲面である。これを繰り返せば全位置で商が得られる。
\end{proof}
